% theme: work
\MajorQuestion{仕事}
\SetRowSpaceQanda

\Instruction{仕事の意味に関する次の問いに答えなさい。}
\QQ{物体に力を加え、その力の向きに物体を移動させたとき、力は物体に\NamedBlank{work}をしたという。}{仕事}
\QQ{\RefBlank{work}の単位を何というか。}{ジュール}
\QQ{あらい水平面上で物体を一定の速さで動かすとき、何と同じ大きさの力を加え続ければよいか。}{摩擦力}
\QQ{力の向きと物体の移動する向きが垂直なとき、\RefBlank{work}の大きさは何Jか。}{0J}
\QQ{物体に力がはたらいていないとき、物体が受ける\RefBlank{work}の大きさは何Jか。}{0J}

\Instruction{仕事の計算方法に関する次の問いに答えなさい。}
\QQ{6Nの力で物体を力の向きに2.4m動かした。この力がした\RefBlank{work}は何Jか。}{14.4J}
\QQ{18Jの\RefBlank{work}で物体を力の向きに3.0m動かした。加えた力は何Nか。}{6N}
\QQ{7.5Nの力で15Jの\RefBlank{work}をした。物体を力の向きに何m動かしたか。}{2m}
\QQ{640gは何kgか。}{0.64kg}
\QQ{125cmは何mか。}{1.25m}

\Instruction{道具を使った仕事に関する次の問いに答えなさい。}
\QQ{道具を使っても使わなくても、\RefBlank{work}の大きさは変わらないという原理を何というか。}{仕事の原理}
\QQ{加える力の向きだけを変え、力の大きさや引く距離を変えない滑車を何というか。}{定滑車}
\QQ{加える力は小さくなるが、引く距離が大きくなる滑車を何というか。}{動滑車}
\QQ{支点のまわりに回転できる棒を使い、小さい力で大きな力を出す道具を\NamedBlank{teko}という。}{てこ}
\QQ{\RefBlank{teko}を使い必要な力を小さくした場合、移動させる距離はどうなるか。}{大きくなる}

\Instruction{仕事率と電気による仕事に関する次の問いに答えなさい。}
\QQ{単位時間あたりにする\RefBlank{work}の大きさを\NamedBlank{power}という。}{仕事率}
\QQ{48Jの\RefBlank{work}を12秒でしたとき、\RefBlank{power}は何Wか。}{4W}
\QQ{\RefBlank{power}が6Wの機械が45Jの\RefBlank{work}をするのにかかる時間は何秒か。}{7.5秒}
\QQ{\RefBlank{power}の単位を何というか。}{ワット}
\QQ{電気による\RefBlank{work}の能率を何というか。}{電力}
